\subsection{哈尔测度的局部定义}

命题 9. --- 设 G 为局部紧群，V 为 G 的开子集，且 \(\mu\) 为 V 上具有如下性质的非零正测度：若 U 是 V 的开子集，且 \(s \in G\) 满足 \(sU \subset V\)，则测度 \(\mu_{U}\)（由 \(\mu\) 在 U 上诱导）在从 U 到 sU 的同胚 \(x \mapsto sx\) 下的像为 \(\mu_{sU}\)。那么，存在唯一一个 G 上的左哈尔测度 \(\alpha\)，它在 V 上诱导出 \(\mu\)。

对于每个 \(s \in G\)，令 \(\mu_{s}\) 为 \(\mu\) 在从 V 到 sV 的同胚 \(x \mapsto sx\) 下的像。将 \(\mu_{s}\) 限制到 \(V \cap sV\)，它是 \(\mu_{s^{-1}V \cap V}\) 在 \(x \mapsto sx\)（定义在 \(s^{-1}V \cap V\) 上）的限制下的像；根据假设，该像为 \(\mu_{V \cap sV}\)。通过平移可知，对于任意 s、t，\(\mu_{s}\) 和 \(\mu_{t}\) 在 \(sV \cap tV\) 上的限制相同。因此根据第三章第 2 节第 1 号命题 1，存在 G 上的测度 \(\alpha\)，使其在每个 sV 上诱导 \(\mu_{s}\)。显然，\(\alpha\) 是 G 上诱导 V 上 \(\mu\) 的唯一左哈尔测度。

推论。--- 设 G、G' 是两个局部紧群，V（分别为 V'）是 G（分别为 G'）中单位元的开邻域，且 \(\varphi\) 是 G' 到 G 的局部同构（GT, III, §1, No.~3, 定义 2），定义在 V' 上，并满足 \(\varphi(V') = V\)。设 \(\alpha'\) 是 G' 上的左哈尔测度，\(\alpha'_{V'}\) 是其在 \(V'\) 上的限制。则 \(\varphi(\alpha_{V'}')\) 是 G 上某个唯一左哈尔测度 \(\alpha\) 在 V 上的限制。

令 \(V_1\) 为 \(e\) 在 \(G\) 中的开邻域，满足 \(V_1V_1^{-1} \subset V\)。令 \(\mu\) 为 \(\varphi(\alpha_{V'}')\) 在 \(V_1\) 上的限制。设 \(U\) 是 \(V_1\) 的开子集，且 \(s \in G\) 满足 \(sU \subset V_1\)。则 \(s \in V_1V_1^{-1} \subset V\)，故 \(s = \varphi(s')\)，其中 \(s' \in V'\)。令 \(x \in U\)。则 \(x = \varphi(x')\)，其中 \(x' \in V'\)，从而 \(sx = \varphi(s')\varphi(x') = \varphi(s'x')\)，因为 \(sx \in sU \subset V\)。由于 \(G'\) 中的左平移保持 \(\alpha'\)，可见 \(V_1\) 和 \(\mu\) 满足命题 9 的条件。令 \(\alpha\) 为 \(G\) 上使 \(\mu\) 在 \(V_1\) 上得到诱导的左哈尔测度。对于每个 \(t \in V\)，存在开邻域 \(W\)，它是 \(e\) 在 \(V_1\) 中的开邻域，使 \(tW \subset V\)。于是 \(\varphi(\alpha_{V'}')\) 在 \(tW\) 上的限制可由 \(\mu\) 在 \(W\) 上的限制通过平移得到，因而是 \(\alpha\) 在 \(tW\) 上的限制。因此 \(\varphi(\alpha_{V'}')\) 是 \(\alpha\) 在 \(V\) 上的限制。

称 \(\alpha\) 是由 \(\alpha'\) 借助局部同构 \(\varphi\) 得到的。

例。--- 第 2 号例 3 中得到的 T 上的哈尔测度，可以通过 R 到 T 的一个局部同构由 R 上的勒贝格测度得到。

